% Standalone solution dossier for affine-Poincare W2 lower semicontinuity
% and the conditional one-sequence CMH recovery implication.
\documentclass[../main.tex]{subfiles}
\begin{document}

% === SOLUTION HEADER (proof provenance) =====================================
% ledger-node : lem:affine-poincare-w2-liminf;
%               cor:cmh-recovery-sequence-suffices
% refines     : lem:affine-poincare-w2-liminf;
%               cor:cmh-recovery-sequence-suffices
%               (modules/kls/41-cmh-normalization.tex)
% bounded_by  : none (both nodes)
% checked_by  : agent
% author      : /root/prove_cmh_recovery
% reviewer    : /root/review_cmh_recovery_w0
% review      : research/reviews/2026-08-27-lem-affine-poincare-w2-liminf-proof-review.md
% date        : 2026-08-27
% ============================================================================

\section*{Solution candidate: affine-Poincar\'e recovery and one-sequence CMH closure}
\label{sec:sol-affine-poincare-w2-liminf}

\paragraph{Scope and certification boundary.}
The first result below proves Lemma~\ref{lem:affine-poincare-w2-liminf}.  In fact, its
lower-semicontinuity assertion is proved along every centered log-concave $W_2$-convergent
sequence, while the existence of compact-target regular approximants uses the published
moment-map result recorded as Theorem~\ref{thm:regular-moment-map-compact-target}.  The second
result proves Corollary~\ref{cor:cmh-recovery-sequence-suffices} under the unresolved
Assumption~\ref{ass:cmh-recovery-envelope}.  No bound on $\CMH$ is proved, and no convergence or
semicontinuity of $\CMH$ is asserted.

\subsection*{1. The covariance form on an affine support}

Let $\mu$ be a centered log-concave probability on $\R^n$, set
\[
  \Sigma=\Cov(\mu),\qquad S=\operatorname{Ran}\Sigma,\qquad d=\dim S,
\]
and, when $d>0$, write $\Sigma_S=\Sigma|_S$.  The affine hull of $\mu$ is the linear space
$S$.  Indeed, if $v\in\ker\Sigma$, then centeredness and
$\Var_\mu\langle v,X\rangle=0$ give $\langle v,X\rangle=0$ almost surely.  Conversely, the
covariance is positive definite on the linear span of the centered support.  Thus
$\Sigma_S\succ0$ for $d>0$.

On the globally Lipschitz functions on $S$ that belong to $L^2(\mu)$, consider the pre-form
\begin{equation}
  \calE^0_{\Sigma,\mu}(f)
  =\int_S\inner{\Sigma_S\nabla_S f}{\nabla_S f}\,\dd\mu.
  \label{eq:sol-cmh-recovery-preform}
\end{equation}
We denote its closure by $\calE_{\Sigma,\mu}$, its closed domain by
$H^1_\Sigma(\mu)$, and define
\begin{equation}
  \CPaff(\mu)
  =\sup_{f\in H^1_\Sigma(\mu)\setminus\R}
       \frac{\Var_\mu f}{\calE_{\Sigma,\mu}(f)}.
  \label{eq:sol-cmh-recovery-cpaff}
\end{equation}
When $d=0$, centeredness gives $\mu=\delta_0$; we set
$H^1_\Sigma(\mu)=L^2(\mu)=\R$ and $\CPaff(\mu)=0$.

\begin{lemma}[Intrinsic and ambient common core]
\label{lem:sol-cmh-recovery-common-core}
For $d>0$, the pre-form \eqref{eq:sol-cmh-recovery-preform} is closable, its closed-form kernel
consists of the constants, and
\[
  \mathscr C_S=\R+C_c^\infty(S)
\]
is a form core.  Moreover, $\mathscr C_S$ is exactly the restriction to $S$ of the one ambient
class
\[
  \mathscr C=\R+C_c^\infty(\R^n),
  \label{eq:sol-cmh-recovery-ambient-core}
\]
and the covariance energy of a restriction is independent of its ambient extension.
\end{lemma}

\begin{proof}
A log-concave probability with affine hull $S$ has a density with respect to Lebesgue measure
on $S$ which is positive and locally bounded above and below on the relative interior of its
convex support.  Suppose that $f_j$ is Lipschitz,
$f_j\to0$ in $L^2(\mu)$, and $\Sigma_S^{1/2}\nabla_Sf_j\to u$ in
$L^2(\mu;S)$.  On every compact subset of the relative interior, the density is equivalent to
Lebesgue measure and $\Sigma_S^{-1/2}$ is bounded.  Hence
$f_j\to0$ and $\nabla_Sf_j\to\Sigma_S^{-1/2}u$ in local Lebesgue $L^2$.
Closedness of distributional differentiation gives $u=0$.  This proves closability.  The same
local argument applied to a zero-energy element of the closure shows that its distributional
gradient vanishes on a connected open convex set, so the element is constant almost
everywhere.

It remains to verify the asserted core rather than assume it.  Start with a Lipschitz
$f\in L^2(\mu)$.  The value truncations $T_Mf=(-M)\vee f\wedge M$ converge to $f$ in $L^2$
and in the energy \eqref{eq:sol-cmh-recovery-preform}, by the Lipschitz chain rule and dominated
convergence.  For bounded $f$, choose $\chi_R\in C_c^\infty(S)$ which equals one on
$B_S(0,R)$, is supported in $B_S(0,2R)$, and satisfies
$|\nabla_S\chi_R|\le c/R$.  Then $\chi_Rf\to f$ in $L^2$, while
\begin{align*}
 \calE^0_{\Sigma,\mu}(\chi_Rf-f)
 &\le 2\int_{S\setminus B_S(0,R)}
          \inner{\Sigma_S\nabla_Sf}{\nabla_Sf}\,\dd\mu \\
 &\quad+\frac{2c^2}{R^2}\norm{f}_\infty^2\norm{\Sigma_S}_\op
 \longrightarrow0.
\end{align*}
Finally, a compactly supported Lipschitz function on the Euclidean space $S$ can be mollified
inside $S$.  Its mollifications converge in value and gradient almost everywhere, and their
gradients are bounded by the original Lipschitz constant.  Absolute continuity of $\mu$ on
$S$ and dominated convergence give convergence in the form norm.  Therefore
$\mathscr C_S$ is a core.

An ambient function in $\mathscr C$ plainly restricts to $\mathscr C_S$.  Conversely, in the
orthogonal splitting $x=s+t\in S\oplus S^\perp$, extend
$\phi\in C_c^\infty(S)$ by
\[
  F(s+t)=\phi(s)\eta(t),
\]
where $\eta\in C_c^\infty(S^\perp)$ equals one near zero.  This is an ambient compactly
supported smooth extension.  If two ambient extensions agree on $S$, the difference of their
gradients on $S$ lies in $S^\perp$.  Since
\begin{equation}
  \Sigma=P_S\Sigma_SP_S,
  \qquad
  \Sigma^+=P_S\Sigma_S^{-1}P_S,
  \label{eq:sol-cmh-recovery-normal-annihilation}
\end{equation}
both $\Sigma$ and its Moore--Penrose inverse annihilate that normal difference.  Thus the
ambient covariance energy is exactly the intrinsic energy and is independent of the extension.
\end{proof}

\subsection*{2. A regular compact-target recovery sequence}

\begin{theorem}[Affine-Poincar\'e $W_2$ lower semicontinuity and regular recovery]
\label{thm:sol-affine-poincare-w2-liminf}
Let $\mu$ be a centered log-concave probability on $\R^n$.
\begin{enumerate}[label=(\roman*),leftmargin=2.2em]
  \item There are centered, full-dimensional, compactly supported log-concave probabilities
        $\mu_k$ in the regular compact-target moment-map class such that
        $\Wtwo(\mu_k,\mu)\to0$ in the original ambient coordinates.
  \item More generally, for every sequence of centered log-concave probabilities $\nu_k$ with
        $\Wtwo(\nu_k,\mu)\to0$,
        \begin{equation}
          \CPaff(\mu)\le\liminf_{k\to\infty}\CPaff(\nu_k).
          \label{eq:sol-affine-poincare-w2-liminf}
        \end{equation}
\end{enumerate}
In particular, the sequence in (i) has the lower-semicontinuity property required by
Lemma~\ref{lem:affine-poincare-w2-liminf}, including when $\mu$ is carried by a proper affine
subspace.  The assertion concerns only covariance Dirichlet forms and makes no claim about
$\CMH(\mu_k)$ or its limit.
\end{theorem}

\begin{proof}
Let $X\sim\mu$ and let $G\sim N(0,I_n)$ be independent.  For $\delta>0$, put
\[
  \lambda_\delta=\mathcal L(X+\sqrt\delta G),
  \qquad q_\delta=\frac{\dd\lambda_\delta}{\dd x}.
\]
The density $q_\delta$ is positive and smooth, and it is log-concave because convolution
preserves log-concavity \cite{BrascampLieb1976}.  The displayed coupling also gives
\begin{equation}
  \Wtwo^2(\lambda_\delta,\mu)\le n\delta,
  \qquad \int x\,\dd\lambda_\delta(x)=0,
  \qquad \Cov(\lambda_\delta)=\Sigma+\delta I_n.
  \label{eq:sol-cmh-recovery-smoothing}
\end{equation}

For $\delta,\eps>0$ and $R<\infty$, define
\begin{equation}
 Z_{\delta,\eps,R}
 =\int_{B(0,R)}q_\delta(x)e^{-\eps|x|^2/2}\,\dd x,
 \qquad
 \dd\widetilde\mu_{\delta,\eps,R}(x)
 =Z_{\delta,\eps,R}^{-1}\one_{B(0,R)}(x)
   q_\delta(x)e^{-\eps|x|^2/2}\,\dd x,
 \label{eq:sol-cmh-recovery-family}
\end{equation}
and let $m_{\delta,\eps,R}$ be its mean.  For fixed $\delta$, as $\eps\downarrow0$ and
$R\uparrow\infty$, the density ratio in \eqref{eq:sol-cmh-recovery-family} relative to
$\lambda_\delta$ tends pointwise to one, its normalizer tends to one, and it is eventually
bounded by two.  Dominated convergence, first with bounded continuous tests and then with
$|x|^2$, gives weak convergence and convergence of second moments.  Hence
\begin{equation}
  \Wtwo(\widetilde\mu_{\delta,\eps,R},\lambda_\delta)\longrightarrow0.
  \label{eq:sol-cmh-recovery-fixed-delta}
\end{equation}

Take $\delta_k=k^{-4}$ and use \eqref{eq:sol-cmh-recovery-fixed-delta} to choose
$0<\eps_k<k^{-1}$ and $R_k>k$ such that
\[
  \Wtwo(\widetilde\mu_{\delta_k,\eps_k,R_k},\lambda_{\delta_k})<k^{-1}.
\]
Set $m_k=m_{\delta_k,\eps_k,R_k}$ and recenter:
\[
  \mu_k=(x\mapsto x-m_k)_\#\widetilde\mu_{\delta_k,\eps_k,R_k}.
\]
Comparison of means under any quadratic-cost coupling and centeredness of $\lambda_{\delta_k}$
give $|m_k|<k^{-1}$.  Therefore the triangle inequality and
\eqref{eq:sol-cmh-recovery-smoothing} yield the ambient estimate
\begin{equation}
  \Wtwo(\mu_k,\mu)<\frac2k+\frac{\sqrt n}{k^2}.
  \label{eq:sol-cmh-recovery-diagonal}
\end{equation}

We verify the regular class before using the words ``regular recovery.''  The support of
$\mu_k$ is the convex body $P_k=B(-m_k,R_k)$, and on its interior
\[
  \dd\mu_k(z)=g_k(z)\,\dd z,
  \qquad
  g_k(z)=Z_{\delta_k,\eps_k,R_k}^{-1}
  q_{\delta_k}(z+m_k)e^{-\eps_k|z+m_k|^2/2}.
\]
The function $g_k$ extends to a positive smooth function on $\R^n$.  It is log-concave,
$\mu_k$ is centered by construction, and it is full-dimensional because $g_k>0$ on the open
ball.  Its barycenter zero lies in $\operatorname{int}P_k$: a positive density on a
full-dimensional convex body cannot have its barycenter on a supporting hyperplane.
The published compact-target theorem
\cite{BermanBerndtsson2013RealMA,Fathi2019SteinMomentMaps}, in the repository form of
Theorem~\ref{thm:regular-moment-map-compact-target}, now supplies a smooth strictly convex
moment potential $\varphi_k$, a diffeomorphism
$\nabla\varphi_k:\R^n\to\operatorname{int}P_k$, and the smooth positive symmetric kernel
\[
 H_k(x)=D^2\varphi_k((\nabla\varphi_k)^{-1}(x))
\]
with weak zero flux,
\begin{equation}
  \Div_{\mu_k}H_k=-x,
  \qquad \E_{\mu_k}H_k=\Sigma_k:=\Cov(\mu_k).
  \label{eq:sol-cmh-recovery-stein}
\end{equation}
Thus (i) is proved.

We next prove the stronger assertion (ii).  Write $\Sigma_k^\nu=\Cov(\nu_k)$.  Quadratic
Wasserstein convergence gives convergence of second moments.  Since the measures are centered,
it follows that
\begin{equation}
  \Sigma_k^\nu\longrightarrow\Sigma
  \label{eq:sol-cmh-recovery-covariance-convergence}
\end{equation}
in every matrix norm.  For completeness, under couplings with
$\E|X_k-X|^2\to0$, Cauchy--Schwarz gives convergence of each
$\E[X_{k,i}X_{k,j}]$ to $\E[X_iX_j]$, because the second moments of $X_k$ remain bounded.

If $d=0$, then $\CPaff(\mu)=0$ by convention and
\eqref{eq:sol-affine-poincare-w2-liminf} is immediate.  We may therefore assume $d>0$ for the
remaining common-core argument.

Fix $F\in\mathscr C$ from \eqref{eq:sol-cmh-recovery-ambient-core}.  Both $F$ and $F^2$ are
bounded and continuous, so weak convergence gives
\begin{equation}
  \Var_{\nu_k}F\longrightarrow\Var_\mu F.
  \label{eq:sol-cmh-recovery-variance-convergence}
\end{equation}
Also, with $h(x)=\inner{\Sigma\nabla F(x)}{\nabla F(x)}$, which is bounded and continuous,
\begin{align}
 &\left|\int\inner{\Sigma_k^\nu\nabla F}{\nabla F}\,\dd\nu_k
       -\int\inner{\Sigma\nabla F}{\nabla F}\,\dd\mu\right| \notag\\
 &\quad\le
 \norm{\Sigma_k^\nu-\Sigma}_\op\norm{\nabla F}_\infty^2
 +\left|\int h\,\dd\nu_k-\int h\,\dd\mu\right|
 \longrightarrow0.
 \label{eq:sol-cmh-recovery-energy-convergence}
\end{align}

Let $L=\liminf_k\CPaff(\nu_k)$.  If $L=\infty$, there is nothing to prove.  Otherwise choose
one subsequence, independent of the test function, along which
$\CPaff(\nu_k)\to L$.  The affine Poincar\'e inequality for $\nu_k$, followed along this fixed
subsequence by \eqref{eq:sol-cmh-recovery-variance-convergence} and
\eqref{eq:sol-cmh-recovery-energy-convergence}, gives for every $F\in\mathscr C$
\[
  \Var_\mu F
  \le L\int\inner{\Sigma\nabla F}{\nabla F}\,\dd\mu.
\]
Lemma~\ref{lem:sol-cmh-recovery-common-core} identifies the restriction of $\mathscr C$ with the
intrinsic core on $S$ and extends the inequality in the closed form norm to all
$H^1_\Sigma(\mu)$.  Hence $\CPaff(\mu)\le L$, proving (ii).  If $\Sigma$ is singular,
\eqref{eq:sol-cmh-recovery-normal-annihilation} removes all normal derivatives; no inverse is
taken in a collapsing direction and no whitening is performed before the limit.
\end{proof}

\subsection*{3. Verification of the regular CMH endpoint}

Although Theorem~\ref{thm:sol-affine-poincare-w2-liminf} itself does not use $\CMH$, the
conditional corollary needs the exact regular endpoint.  We record why the compact-target
objects above have the closed-form realization required by Definition~\ref{def:cmh}.

Here a regular recovery sequence means a sequence of centered, full-dimensional,
compact-target moment-map laws carrying the canonical kernel and weak zero-flux convention of
Theorem~\ref{thm:regular-moment-map-compact-target}.  For one such law $\nu$, with target $P$,
density $g$, covariance $\Sigma_\nu\succ0$, and canonical kernel $H$, let
\[
  \mathscr D=\{F|_P:F\in\R+C_c^\infty(\R^n)\},
  \qquad
  \calE_H^0(f)=\int_P\inner{H\nabla f}{\nabla f}\,\dd\nu.
\]
This core is dense in $L^2(\nu)$ because it contains the restrictions of
$C_c^\infty(\operatorname{int}P)$ and the convex boundary is null.  The form is finite there
because $\E_\nu H=\Sigma_\nu$.  It is closable: if
$f_j\to0$ in $L^2(\nu)$ and $H^{1/2}\nabla f_j\to u$ in $L^2(\nu;\R^n)$, then on each compact
subset of $\operatorname{int}P$, the density is bounded above and below and $H^{\pm1/2}$ are
bounded.  Thus $f_j\to0$ and $\nabla f_j\to H^{-1/2}u$ in local Lebesgue $L^2$; closedness of
distributional differentiation forces $u=0$.  The same argument shows that a zero-energy
element of the closure is constant, because $\operatorname{int}P$ is connected.  Finally, the
global weak Stein identity in \eqref{eq:sol-cmh-recovery-stein} identifies the Friedrichs form
operator with the closed Stein generator used in Definition~\ref{def:cmh}; there is no hidden
boundary distribution or unnamed maximal-domain convention.  Consequently the already-certified
Theorem~\ref{thm:cmh-implies-affine-poincare} applies and gives
\begin{equation}
  \CPaff(\nu)\le\CMH(\nu)
  \label{eq:sol-cmh-recovery-regular-endpoint}
\end{equation}
for every member of a regular recovery sequence.

\begin{corollary}[One bounded CMH recovery sequence suffices]
\label{cor:sol-cmh-recovery-sequence-suffices}
Assume Assumption~\ref{ass:cmh-recovery-envelope}: there is a universal $C$ such that, for every
centered log-concave $\mu$, at least one centered regular recovery sequence $\mu_k$ satisfies
\[
  \Wtwo(\mu_k,\mu)\to0,
  \qquad
  \liminf_{k\to\infty}\CMH(\mu_k)\le C.
\]
Then every centered log-concave probability, including one carried by a proper affine
subspace, satisfies
\[
  \CPaff(\mu)\le C.
\]
Thus the recovery-envelope assumption implies Conjecture~\ref{conj:kls} with the same constant.
\end{corollary}

\begin{proof}
Apply the general lower-semicontinuity assertion
\eqref{eq:sol-affine-poincare-w2-liminf} to the one sequence supplied by the assumption, and
apply \eqref{eq:sol-cmh-recovery-regular-endpoint} to each of its regular members.  Pointwise
comparison of the two sequences of constants gives
\[
  \CPaff(\mu)
  \le\liminf_k\CPaff(\mu_k)
  \le\liminf_k\CMH(\mu_k)
  \le C.
\]
All limits are taken in the original ambient coordinates.  The singular-support convention is
therefore exactly that of Lemma~\ref{lem:sol-cmh-recovery-common-core}, and the constant does not
change.  A universal bound on $\CPaff$ is the affine-invariant Poincar\'e formulation of KLS.
\end{proof}

\paragraph{Hypotheses and conditional status.}
The lemma uses centeredness, log-concavity, finite-dimensionality, and the published regular
compact-target moment-map theorem.  Centeredness identifies the affine hull with
$\operatorname{Ran}\Sigma$ and is preserved by the construction.  Log-concavity supplies the intrinsic
density and is preserved by smoothing, tilt, convex truncation, translation, and $W_2$ limits.
No isotropy, spectral gap, smoothness of the limiting law, or full-dimensionality of the limit
is assumed.  The corollary additionally uses Definition~\ref{def:cmh}, the certified endpoint
Theorem~\ref{thm:cmh-implies-affine-poincare}, and the unresolved
Assumption~\ref{ass:cmh-recovery-envelope}.  The lemma is unconditional relative to its
published imported moment-map input; the corollary remains conditional precisely on that open
recovery-envelope assumption.

\paragraph{Obstructions respected.}
Neither ledger node has a \texttt{bounded\_by} edge.  The proof uses no localization occupation
estimate, fixed cut, projection-only test, trace upgrade, or evolving isoperimetric competitor.
Affine-support collapse is handled solely by ambient $W_2$ convergence and the intrinsic closed
covariance form.  No canonical moment-map kernel is transported through a noninvertible map,
and no lower semicontinuity of $\CMH$ is claimed.

\paragraph{Unclosed step and deferred ledger artifact.}
The analytic proof of the lemma has no unclosed step.  The only unclosed mathematical premise in
the corollary is Assumption~\ref{ass:cmh-recovery-envelope}.  Since this dossier has
\texttt{checked\_by: none}, it has no ledger value.  The future shared
\texttt{solution: solutions/lem-affine-poincare-w2-liminf.tex} is only a deferred artifact
candidate pending independent review; no ledger delta is applicable now.

\end{document}
